% 原书第 195—198 页；完整章首导言与第 13.1 节。
\chapter{整群随机试验中的数据自适应估计}
\label{ch:data-adaptive-cluster-randomized-trials}
\chapterauthors{Laura B. Balzer、Mark J. van der Laan 与 Maya L. Petersen}
\chapteraffiliations{%
L. B. Balzer（通讯作者）：美国马萨诸塞大学阿默斯特分校公共卫生与健康科学学院生物统计学与流行病学系，416 Arnold House, 715 North Pleasant St, Amherst, MA 01003, USA；电子邮箱：\href{mailto:lbalzer@umass.edu}{lbalzer@umass.edu}。\\
M. J. van der Laan：美国加利福尼亚大学伯克利分校生物统计学部及统计学系，101 Haviland Hall, \#7358, Berkeley, CA 94720, USA；电子邮箱：\href{mailto:laan@berkeley.edu}{laan@berkeley.edu}。\\
M. L. Petersen：美国加利福尼亚大学伯克利分校流行病学部及生物统计学部，101 Haviland Hall, \#7358, Berkeley, CA 94720, USA；电子邮箱：\href{mailto:mayaliv@berkeley.edu}{mayaliv@berkeley.edu}。\\
本章 DOI：\url{https://doi.org/10.1007/978-3-319-65304-4_13}。}

在随机试验中，分析时调整已测量的协变量，可以减小方差、提高统计功效。为避免产生误导性的推断，必须预先指定分析方案。然而，我们往往事先并不清楚，应当在分析中纳入哪些基线协变量，甚至不清楚是否应纳入任何协变量。由此带来一个重要挑战：我们需要从数据中学习，以提高精度；同时，又必须以预先指定的严格方式学习，保持统计推断的有效性。这一挑战在整群随机试验（cluster randomized trial，CRT）中尤其突出，因为这类试验的独立单位，如社区、诊所或学校，数量往往有限，而可能用于调整的变量却很多。

\begin{sourcebox}
本章讨论一种严格的方法，以数据自适应方式选择调整集，使分析的效率最大化。具体而言，我们使用交叉验证，从预先指定的算法库中选择使估计方差最小的候选靶向最小损失估计量（targeted minimum loss-based estimator，TMLE）。为进一步提高精度，我们还提出一种协同方法，用来估计已知的暴露机制。小样本模拟表明，这种方法有望使研究功效最大化，同时保持置信区间的名义覆盖概率。我们说明如何针对科学问题，即研究样本或目标总体中的干预效应，以及研究设计，即是否采用配对匹配，来调整这一方法。本章改编自 Balzer et al.（2016b）。
\end{sourcebox}

随机试验旨在评价某项干预对感兴趣结局的影响。在这种情形下，接受处理的单位与对照单位之间的平均结局之差，是干预效应的一个简单无偏估计量。在分析时调整已测量的协变量，可以显著减小估计量的方差，从而提高研究功效，例如 Fisher（1932）、Cochran（1957）、Cox and McCullagh（1982）、Tsiatis et al.（2008）以及 Moore and van der Laan（2009b）。然而，关于随机试验中应如何调整、何时调整的建议，却一直相互矛盾（ICH Harmonised Tripartite Guideline，1998；Pocock et al.，2002；Hayes and Moulton，2009；Austin et al.，2010；Kahn et al.，2014；Campbell，2014；European Medicines Agency，2015）。这些建议似乎取决于研究设计、随机化单位、应用领域和样本量。因此，许多研究者仍不知道应如何调整基线协变量，甚至不知道是否应作调整。

令 $n$ 为研究单位，如社区或学校的数量。考虑一项试验：将 $n/2$ 个单位随机分配至处理组，其余单位分配至对照组。关于这一情形下的局部有效估计，已有丰富的文献，例如 Tsiatis et al.（2008）、Zhang et al.（2008）、Rubin and van der Laan（2008）、Moore and van der Laan（2009b）以及 Shen et al.（2014）。例如，参数回归可以给出干预效应的一个无偏、且精度更高的估计。其基本做法是，根据一个工作模型，将结局对暴露和协变量作回归。按照 Rosenblum and van der Laan（2010b）的用法，我们使用“工作”一词，强调回归函数不必设定正确，而且往往确实存在设定错误。工作模型可以包含交互项，也可以是线性或非线性的。随后，用估计的系数预测所有单位分别在处理和对照条件下的结局。对预测结局分别取平均，再计算二者之差或之比，就得到干预效应的估计。在观察性研究中，这一算法有时称为参数 g 计算（parametric g-computation；Robins，1986）。

对于连续结局以及不含交互项的线性工作模型，这个方法称为协方差分析（analysis of covariance，ANCOVA；Cochran，1957），暴露的系数就等于干预效应的估计值。对于二值结局，Moore and van der Laan（2009b）详细说明了，通过 logistic 回归作调整，在绝对或相对尺度上估计处理效应，即风险差、风险比或优势比时，可能获得的精度改善。此外，他们证明了，在这种情形下，参数极大似然估计等价于 TMLE。因此，即使结局回归的工作模型存在设定错误，TMLE 的双重稳健性、渐近线性等渐近性质仍然成立。而且，这种方法具有局部有效性：如果工作模型设定正确，TMLE 就会在一大类估计量中达到可能的最小方差。Rosenblum and van der Laan（2010b）将这些结果推广到一大类一般线性模型。事实上，参数极大似然估计量（maximum likelihood estimator，MLE）和 TMLE，都可以视为 Scharfstein et al.（1999）的双重稳健估计量，以及 Tsiatis et al.（2008）和 Zhang et al.（2008）的半参数方法的特例。关于这些估计方法近期的详细综述，可参阅 Colantuoni and Rosenblum（2015）。

现在考虑配对匹配试验：在 $n/2$ 个匹配对内随机分配干预。针对这一情形提出的估计策略较为有限。事实上，人们认为配对匹配试验存在“分析上的局限”，以致一些研究者回避这种设计（Klar and Donner，1997；Imbens，2011；Campbell，2014）。与完全随机试验一样，未经调整的特定处理下均值之差，可以给出干预效应的无偏估计，但效率不高。为了在分析中纳入协变量、提高功效，Hayes and Moulton（2009）建议，将结局对协变量作回归，但不对暴露作回归，再在匹配对内比较观测结局与预测结局。另一种选择是采用 TMLE，它在配对匹配试验中可以提供无偏、局部有效的方法（van der Laan et al.，2013a；Balzer et al.，2015，2016c）。具体而言，可以按照完全随机试验的方式实施算法：（1）拟合给定暴露和协变量时结局均值的工作模型；（2）预测所有单位分别在处理和对照条件下的结局；（3）在相关尺度上比较这些预测结局的均值。不过，推断必须遵循配对匹配方案。

这两种设计都面临一个共同挑战：选择分析中应纳入的协变量。在实施干预前，往往已经测量了许多变量，事先指定一个合适的工作模型并不容易。在完全随机试验中，如果（a）协变量能够预测结局，而且（b）这些协变量在处理组之间不平衡，那么调整协变量就能提高精度，例如 Moore et al.（2011）。当样本量趋于无穷时，可以保证平衡；但在实际中，完全平衡却很少见。同样，在配对匹配试验中，如果匹配后仍然存在预后协变量的不平衡，协变量调整就能提高精度。

有限的样本量给协变量选择带来另一重挑战。近期一项随机临床试验综述报告，参与者数量的中位数为 $58$，四分位数区间为 $27$—$161$（Califf et al.，2012）。同样，近期一项 CRT 综述报告，簇数的中位数为 $31$，四分位数区间为 $13$—$60$（Selvaraj and Prasad，2013）。在小规模试验中，调整过多协变量可能导致过拟合，并使第一类错误率升高，例如 Moore et al.（2011）、Shen et al.（2014）以及 Balzer et al.（2015）。最后，临时决定调整集也会引发担忧：研究者可能进行“撒网式搜索”，寻找能使功效最大的协变量，从而再次面临第一类错误率升高的风险，例如 Pocock et al.（2002）、Tsiatis et al.（2008）、Olken（2015）以及 Rose（2015）。

总之，在随机试验中调整协变量，可以显著提高精度，进而提高统计功效。为避免误导性的统计推断，必须事先指定工作模型，包括调整变量。实际中，样本量往往限制了调整集的大小，而在试验结束前，最佳调整集又并不明确。本章应用经验效率最大化原则（empirical efficiency maximization），以数据自适应方式，从预先指定的算法库中选择使方差最小的候选 TMLE，从而使分析精度最大化（Rubin and van der Laan，2008）。我们修改这一策略，使其适用于配对匹配试验；还协同估计暴露机制，以进一步提高精度（van der Laan and Gruber，2010）。我们也将估计与推断的结果，推广至总体平均处理效应和样本平均处理效应（Neyman，1923；Balzer et al.，2016c）。有限样本模拟展示了，在独立单位数量有限时，这些方法的实际表现；这正是早期临床试验和 CRT 中常见的情形。

\section{研究实例与因果参数}
\label{sec:ch13-motivating-example}

作为研究实例，考虑一项社区随机试验，估计立即开展抗逆转录病毒治疗（ART）对 HIV 发生率的影响。假设试验在 $30$ 个社区开展，并在基线测量了大量协变量。再假设，用其中一部分特征形成了 $15$ 个最佳社区匹配对（Balzer et al.，2015）。主要结局是相关时间段内 HIV 的累积发生率。某个研究社区的观测数据可以记为 $O=(W,A,Y)$，其中，$W$ 表示基线协变量向量，$A$ 表示干预分配，$Y$ 表示结局。具体而言，$W$ 包括地区、HIV 流行率、男性包皮环切覆盖率以及社区层面的 HIV RNA 病毒载量；$A$ 为二值指示变量，若社区被随机分配至处理组则等于 $1$，若被随机分配至对照组则等于 $0$；$Y$ 为估计的 HIV 累积发生率。我们关注总体平均处理效应（population average treatment effect，PATE）的估计与推断：
\begin{equation}
 \mathrm{PATE}=\mathbb E[Y_1-Y_0],
 \tag{13.1}
\end{equation}
以及样本平均处理效应（sample average treatment effect，SATE）：
\begin{equation}
 \mathrm{SATE}=\frac1n\sum_{i=1}^n[Y_{1,i}-Y_{0,i}].
 \tag{13.2}
\end{equation}
关于这些因果参数的详细讨论，可参阅第 12 章。

